% =========================================================
%  Paper: Implementacion de un Perceptron desde Cero
%  Plantilla RATE — Revista Argentina de Trabajos Estudiantiles
%  Formato dos columnas, A4, estilo IEEE
%  Compatible con pdfLaTeX (TeXpage)
% =========================================================
\documentclass[10pt,twocolumn,a4paper]{article}

% ----------------------------------------------------------
% Paquetes
% ----------------------------------------------------------
\usepackage[T1]{fontenc}
\usepackage[utf8]{inputenc}
\usepackage[spanish,provide=*]{babel}
\usepackage{times}
\usepackage{geometry}
\usepackage{graphicx}
\usepackage{array}
\usepackage{caption}
\usepackage{amsmath}
\usepackage{amssymb}
\usepackage{lettrine}
\usepackage{microtype}
\usepackage{setspace}
\usepackage{titlesec}
\usepackage{fancyhdr}
\usepackage{lastpage}
\usepackage{footnote}
\usepackage{dblfloatfix}
\usepackage{url}
\usepackage{listings}
\usepackage{xcolor}

% ----------------------------------------------------------
% Geometria de pagina
% ----------------------------------------------------------
\geometry{
  a4paper,
  top=19mm,
  bottom=30mm,
  left=13mm,
  right=13mm,
  columnsep=4mm
}

% ----------------------------------------------------------
% Configuracion listings
% ----------------------------------------------------------
\lstset{
  language=Python,
  basicstyle=\ttfamily\fontsize{7}{8}\selectfont,
  keywordstyle=\bfseries,
  commentstyle=\itshape,
  stringstyle=\ttfamily,
  breaklines=true,
  frame=single,
  numbers=left,
  numberstyle=\tiny,
  numbersep=4pt,
  showstringspaces=false,
  tabsize=2,
  captionpos=b
}

% ----------------------------------------------------------
% Tipografia y espaciado
% ----------------------------------------------------------
\setlength{\parindent}{3.5mm}
\setlength{\parskip}{0pt}

% ----------------------------------------------------------
% Encabezado / pie
% ----------------------------------------------------------
\pagestyle{fancy}
\fancyhf{}
\fancyhead[R]{\thepage}
\renewcommand{\headrulewidth}{0pt}

% ----------------------------------------------------------
% Formato de secciones
% ----------------------------------------------------------
\titleformat{\section}
  {\normalsize\bfseries\centering\scshape}
  {\Roman{section}.}{0.5em}{}
\titlespacing{\section}{0pt}{6pt}{4pt}

\titleformat{\subsection}
  {\normalsize\itshape}
  {\Alph{subsection}.}{0.5em}{}
\titlespacing{\subsection}{0pt}{4pt}{2pt}

% ----------------------------------------------------------
% Captions
% ----------------------------------------------------------
\captionsetup[figure]{
  labelsep=period,
  font={small},
  justification=centering,
  position=below
}
\captionsetup[table]{
  labelsep=newline,
  font={small,sc},
  justification=centering,
  position=above,
  name=TABLA
}

% ----------------------------------------------------------
% Titulo en una sola columna
% ----------------------------------------------------------
\makeatletter
\renewcommand{\maketitle}{%
  \twocolumn[{%
    \begin{center}
      {\fontsize{24}{28}\selectfont\@title\par}
      \vspace{4mm}
      {\normalsize\@author\par}
    \end{center}
    \vspace{4mm}
  }]%
  \footnotetext{$^*$ Trabajo presentado como actividad academica en la
    asignatura de Inteligencia Artificial, Universidad Andina del
    Cusco.}%
}
\makeatother

% ----------------------------------------------------------
% Metadatos
% ----------------------------------------------------------
\title{Implementacion de un Perceptron desde Cero:\
       Clasificacion Binaria con Multiples Funciones de Activacion\
       y Validacion Multilenguaje}

\author{%
  Honorio Morales Ttito$^*$\hfill Escuela Profesional de Ingenieria de Sistemas \\
  \textit{Universidad Andina del Cusco, Cusco, Peru}
}

\date{}

% =========================================================
\begin{document}
% =========================================================

\maketitle
\thispagestyle{fancy}

% ----------------------------------------------------------
% Resumen / Abstract
% ----------------------------------------------------------
\noindent\textbf{\textit{Resumen}}---Este trabajo presenta la
implementacion de un perceptron desde cero para clasificacion
binaria, sin usar librerias de aprendizaje automatico como
scikit-learn o TensorFlow/Keras. El sistema integra seis funciones de
activacion (lineal, escalon, sigmoidal, ReLU, softmax y tangente
hiperbolica) y se evalua en seis casos: AND, OR, spam, clima, fraude
y riesgo academico. Ademas de la version original en Python, se
organizo una arquitectura multilenguaje con implementaciones en Java,
JavaScript y C\#, junto con una pagina web estatica para navegacion
del proyecto. En Python se obtuvo convergencia perfecta para AND y OR
(accuracy = 1.0), mientras que en problemas sinteticos se observo
accuracy entre 0.62 y 0.96.

\smallskip
\noindent\textbf{\textit{Abstract}}---This paper presents a
from-scratch perceptron implementation for binary classification,
without using machine learning libraries such as scikit-learn or
TensorFlow/Keras. The system includes six activation functions
(linear, step, sigmoid, ReLU, softmax, and hyperbolic tangent) and is
evaluated on six tasks: AND, OR, spam, weather, fraud, and academic
risk. Beyond the original Python version, a multi-language
architecture was organized with Java, JavaScript, and C\#
implementations, plus a static web page to navigate project contents.
Python achieved perfect convergence for AND and OR (accuracy = 1.0),
while synthetic tasks showed accuracy between 0.62 and 0.96.

% ----------------------------------------------------------
\section{Introduccion}
% ----------------------------------------------------------

\lettrine[lines=2]{E}{}l perceptron, propuesto por Rosenblatt en
1958~[1], representa una base teorica y practica para comprender el
aprendizaje supervisado en redes neuronales. Aunque existen
frameworks modernos de alto nivel, la implementacion desde cero
permite analizar con detalle la actualizacion de pesos, el rol de la
funcion de activacion y las limitaciones de la separabilidad lineal.

El proyecto se inicio con una implementacion funcional en Python y fue
extendido a una arquitectura multilenguaje para cumplir una entrega
academica completa. Actualmente el repositorio contiene:
\begin{itemize}
  \item nucleo del perceptron y activaciones,
  \item seis casos de clasificacion binaria,
  \item exportacion de resultados y graficas,
  \item implementaciones equivalentes en Java, JavaScript y C\#,
  \item pagina web estatica para documentacion y navegacion.
\end{itemize}

El repositorio del proyecto se encuentra en:
\begin{center}
\small\url{https://github.com/Honorio-Morales/perceptrondiez}
\end{center}

% ----------------------------------------------------------
\section{Fundamento Teorico}
% ----------------------------------------------------------

\subsection{Modelo del Perceptron}

Un perceptron monocapa calcula su salida mediante:
\begin{equation}
  \hat{y} = f\!\left(\sum_{j=1}^{n} w_j x_j + b\right) = f(\mathbf{w}^T\mathbf{x} + b)
\end{equation}

\noindent donde $\mathbf{x} \in \mathbb{R}^n$ es el vector de
entradas, $\mathbf{w} \in \mathbb{R}^n$ el vector de pesos y $b$ el
sesgo. La regla de aprendizaje por muestra es:
\begin{equation}
  \Delta w_j = \eta \cdot (y - \hat{y}) \cdot f'(z) \cdot x_j
\end{equation}

\noindent con $\eta$ como tasa de aprendizaje y
$z = \mathbf{w}^T\mathbf{x} + b$. Para la funcion escalon se emplea
la forma clasica de perceptron tomando $f'(z)=1$.

\subsection{Funciones de Activacion}

Las funciones implementadas y su rango de salida se resumen en la
Tabla~I.

\begin{table}[t]
\caption{Funciones de Activacion Implementadas}
\label{tab:activaciones}
\centering
\small
\begin{tabular}{p{1.5cm}|p{3.2cm}|p{1.3cm}}
\hline
\textbf{Funcion} & \textbf{Expresion} & \textbf{Rango} \\
\hline
Lineal     & $f(z) = z$ & $(-\infty,\infty)$ \\
\hline
Escalon    & $f(z) = \begin{cases}1 & z \ge 0\\0 & z < 0\end{cases}$ & $\{0,1\}$ \\
\hline
Sigmoidal  & $f(z) = \dfrac{1}{1+e^{-z}}$ & $(0,1)$ \\
\hline
ReLU       & $f(z) = \max(0, z)$ & $[0,\infty)$ \\
\hline
Softmax    & $f(z_i) = \dfrac{e^{z_i}}{\sum_j e^{z_j}}$ & $(0,1)$ \\
\hline
Tanh       & $f(z) = \tanh(z)$ & $(-1,1)$ \\
\hline
\end{tabular}
\end{table}

% ----------------------------------------------------------
\section{Metodologia}
% ----------------------------------------------------------

\subsection{Arquitectura del Software}

La solucion se organiza por capas y por lenguaje:
\begin{itemize}
  \item \textbf{Python:} modulo principal de referencia
  (\texttt{main.py}, \texttt{nucleo/}, \texttt{casos/}).
  \item \textbf{Implementaciones:}
  \texttt{implementaciones/java/},
  \texttt{implementaciones/javascript/},
  \texttt{implementaciones/csharp/}.
  \item \textbf{Presentacion:} pagina estatica
  \texttt{web/index.html} para mostrar alcance del proyecto.
\end{itemize}

Cada implementacion mantiene la restriccion de epocas mayor o igual a
10 y ejecuta seis casos de clasificacion binaria.

\subsection{Conjunto de Datos}

Los casos AND y OR usan tablas de verdad clasicas. Los cuatro casos
restantes se generan de manera sintetica con semillas fijas para
reproducibilidad:
\begin{itemize}
  \item \textbf{Spam (100 muestras):} longitud de asunto, numero de
  enlaces, presencia de oferta y remitente desconocido.
  \item \textbf{Clima (100 muestras):} temperatura, humedad y presion
  normalizadas.
  \item \textbf{Fraude (100 muestras):} monto, hora, distancia e
  intentos fallidos.
  \item \textbf{Riesgo academico (100 muestras):} asistencia,
  promedio, entregas pendientes y horas de estudio.
\end{itemize}

\subsection{Metricas de Evaluacion}

Se utilizaron las metricas Accuracy, Precision y Recall:
\begin{equation}
  \text{Accuracy} = \frac{VP + VN}{VP + VN + FP + FN}
\end{equation}

\begin{equation}
  \text{Precision} = \frac{VP}{VP + FP}
  \qquad
  \text{Recall} = \frac{VP}{VP + FN}
\end{equation}

Adicionalmente, para analizar estabilidad de aprendizaje se monitoreo
el error cuadratico medio (ECM) por epoca:
\begin{equation}
  \text{ECM}_t = \frac{1}{N}\sum_{i=1}^{N}(y_i-\hat{y}_i)^2
\end{equation}

\subsection{Flujo Experimental}

El flujo seguido para cada lenguaje fue uniforme:
\begin{enumerate}
  \item inicializacion de pesos en cero y sesgo nulo,
  \item seleccion de funcion de activacion segun el caso,
  \item entrenamiento por epocas con actualizacion muestra a muestra,
  \item evaluacion con Accuracy, Precision y Recall,
  \item consolidacion de resultados para comparacion cruzada.
\end{enumerate}

Este diseno permitio evitar sesgos por diferencias de pipeline y
centrar el analisis en comportamiento numerico y convergencia.

\subsection{Complejidad Computacional}

Para un perceptron monocapa con $N$ muestras, $d$ caracteristicas y
$E$ epocas, el costo de entrenamiento es:
\begin{equation}
  \mathcal{O}(E\cdot N\cdot d)
\end{equation}

Como los cuatro casos sinteticos usan $N=100$ y $d\in\{3,4\}$, los
tiempos de ejecucion son bajos en todos los lenguajes analizados,
facilitando pruebas repetidas y ajustes de hiperparametros.

% ----------------------------------------------------------
\section{Resultados}
% ----------------------------------------------------------

\subsection{Resultados de Referencia en Python}

La Tabla~II presenta los resultados oficiales obtenidos en la
implementacion Python, que se utiliza como referencia del proyecto.

\begin{table}[t]
\caption{Resultados de los Seis Casos en Python}
\label{tab:resultados_python}
\centering
\small
\begin{tabular}{p{2.1cm}|c|c|c|c}
\hline
\textbf{Caso} & \textbf{Activ.} & \textbf{Epocas} & \textbf{Acc.} & \textbf{Prec.} \\
\hline
AND Logico       & escalon   & 15 & 1.0000 & 1.0000 \\
OR Logico        & escalon   & 15 & 1.0000 & 1.0000 \\
Spam             & sigmoidal & 50 & 0.8300 & 0.8333 \\
Clima            & tanh      & 50 & 0.7000 & 0.6296 \\
Fraude           & relu      & 50 & 0.6200 & 0.0000 \\
Riesgo Academico & sigmoidal & 50 & 0.9600 & 0.9649 \\
\hline
\end{tabular}
\end{table}

\subsection{Comparacion Multilenguaje}

La Tabla~III compara accuracy por lenguaje para los entornos
actualmente ejecutables (Python, Java y JavaScript).

\begin{table}[t]
\caption{Comparacion de Accuracy por Lenguaje}
\label{tab:multi}
\centering
\small
\begin{tabular}{p{1.8cm}|c|c|c}
\hline
\textbf{Caso} & \textbf{Python} & \textbf{Java} & \textbf{JS} \\
\hline
AND    & 1.0000 & 1.0000 & 1.0000 \\
OR     & 1.0000 & 1.0000 & 1.0000 \\
Spam   & 0.8300 & 0.7700 & 0.8400 \\
Clima  & 0.7000 & 0.6300 & 0.6300 \\
Fraude & 0.6200 & 0.6600 & 0.6600 \\
Riesgo & 0.9600 & 0.8900 & 0.9000 \\
\hline
\end{tabular}
\end{table}

En C\# los seis casos ya estan implementados, pero en el entorno de
prueba no fue posible ejecutar validacion automatica por ausencia del
comando \texttt{dotnet}.

\subsection{Precision y Recall por Lenguaje}

La Tabla~IV agrega precision y recall para Java y JavaScript,
permitiendo interpretar mejor los falsos positivos y falsos
negativos, especialmente en fraude y clima.

\begin{table}[t]
\caption{Metricas Detalladas en Java y JavaScript}
\label{tab:detalladas_multi}
\centering
\small
\begin{tabular}{p{1.5cm}|c|c|c|c}
\hline
\textbf{Caso} & \textbf{J-Prec} & \textbf{J-Rec} & \textbf{JS-Prec} & \textbf{JS-Rec} \\
\hline
AND    & 1.0000 & 1.0000 & 1.0000 & 1.0000 \\
OR     & 1.0000 & 1.0000 & 1.0000 & 1.0000 \\
Spam   & 0.6667 & 0.9545 & 0.7681 & 1.0000 \\
Clima  & 0.5647 & 1.0000 & 0.5843 & 1.0000 \\
Fraude & 0.0000 & 0.0000 & 0.0000 & 0.0000 \\
Riesgo & 0.9487 & 0.8043 & 0.9348 & 0.8600 \\
\hline
\end{tabular}
\end{table}

\subsection{Analisis por Caso}

\textbf{AND y OR.} Ambos problemas convergen de forma exacta, lo cual
era esperable por su separabilidad lineal. Este resultado funciona
como prueba de sanidad para validar la implementacion base del
algoritmo en cada lenguaje.

\textbf{Spam.} La activacion sigmoidal mantiene un desempeno robusto
en los tres entornos ejecutados. El recall alto en Java y JavaScript
sugiere buena recuperacion de la clase positiva, aunque con diferente
nivel de precision.

\textbf{Clima.} Se observa sensibilidad al umbral con tanh. Aunque el
recall es alto en Java/JS, la precision moderada indica que la
frontera lineal propuesta no separa limpiamente todas las regiones del
espacio de entrada.

\textbf{Fraude.} Es el escenario mas desafiante. Precision y recall
nulos en Java/JS, y precision nula en Python, muestran que la
configuracion actual (activacion relu + una sola neurona) no captura
adecuadamente el patron sintetico.

\textbf{Riesgo academico.} Se mantiene como el mejor caso sintetico,
con metricas altas y comportamiento estable entre implementaciones.

% ----------------------------------------------------------
\section{Discusion}
% ----------------------------------------------------------

\lettrine[lines=2]{L}{}os resultados confirman que el perceptron
resuelve de forma exacta problemas linealmente separables (AND y OR)
en todos los lenguajes evaluados. En los casos sinteticos se observan
diferencias moderadas entre implementaciones, atribuibles a
particularidades numericas de cada runtime y a variaciones en los
generadores pseudoaleatorios.

La activacion sigmoidal mantiene el mejor comportamiento global para
clasificacion binaria en este proyecto. En contraste, el caso de
fraude conserva baja precision, lo que sugiere combinacion de baja
separabilidad lineal y sensibilidad de ReLU en un modelo monocapa.
Esto es consistente con las limitaciones teoricas del perceptron como
clasificador lineal.

Tambien se debe considerar que la comparacion multilenguaje no es una
prueba de equivalencia bit a bit: cada runtime usa generadores
aleatorios y operaciones de punto flotante con pequenas diferencias
internas. Por ello, se priorizo consistencia metodologica sobre
identidad numerica exacta.

\subsection{Amenazas a la Validez}

\begin{itemize}
  \item \textbf{Validez interna:} el entrenamiento y evaluacion se
  realizan sobre el mismo conjunto en varios casos, lo que puede
  sobreestimar desempeno.
  \item \textbf{Validez externa:} los datos de spam, clima, fraude y
  riesgo son sinteticos; no representan toda la complejidad del mundo
  real.
  \item \textbf{Validez de constructo:} softmax se adapta a un
  contexto binario simplificado en perceptron monocapa; en escenarios
  multiclase se requeriria formulacion especifica.
\end{itemize}

% ----------------------------------------------------------
\section{Conclusiones}
% ----------------------------------------------------------

\begin{enumerate}
  \item Se implemento un perceptron desde cero con seis funciones de
  activacion y restriccion de epocas mayor o igual a 10.
  \item Se logro cobertura funcional de seis casos de clasificacion
  binaria en Python, Java y JavaScript, y se completo la base
  equivalente en C\#.
  \item La arquitectura actual del repositorio facilita desarrollo
  colaborativo, comparacion de resultados y trazabilidad academica.
  \item La pagina web estatica complementa la entrega al presentar de
  manera clara el alcance y las implementaciones del proyecto.
\end{enumerate}

Como trabajo futuro se propone incorporar particion entrenamiento/prueba
uniforme en todos los lenguajes, calibracion de umbrales y extension
hacia perceptron multicapa con retropropagacion.

% ----------------------------------------------------------
\section*{Anexo Tecnico}
% ----------------------------------------------------------

Para fines de trazabilidad, el siguiente pseudocodigo resume el ciclo
de entrenamiento utilizado en las implementaciones:

\begin{lstlisting}[language=Python,caption={Pseudocodigo del entrenamiento del perceptron}]
for epoca in range(epocas):
  for i in range(n_muestras):
    z = dot(X[i], pesos) + sesgo
    y_hat = activacion(z)
    error = y[i] - y_hat

    if activacion == "escalon":
      ajuste = lr * error
    else:
      ajuste = lr * error * derivada(z)

    pesos += ajuste * X[i]
    sesgo += ajuste
\end{lstlisting}

Esta formulacion se mantuvo conceptualmente constante en Python, Java,
JavaScript y C\#, permitiendo comparar comportamiento del modelo bajo
un mismo criterio de aprendizaje.

% ----------------------------------------------------------
\section*{Repositorio de Codigo}
% ----------------------------------------------------------

\begin{center}
\small\url{https://github.com/Honorio-Morales/perceptrondiez}
\end{center}

% ----------------------------------------------------------
\section*{Comandos de Reproduccion}
% ----------------------------------------------------------

\begin{lstlisting}[language=bash, numbers=none, frame=single,
  basicstyle=\ttfamily\fontsize{8}{9}\selectfont]
# Python
/home/honorio/IA/perceptron/venv/bin/python main.py

# Java
javac implementaciones/java/src/*.java
java -cp implementaciones/java/src Main

# JavaScript
node implementaciones/javascript/main.js
\end{lstlisting}

% ----------------------------------------------------------
\section*{Reconocimientos}
% ----------------------------------------------------------

Se agradece a la Escuela Profesional de Ingenieria de Sistemas de la
Universidad Andina del Cusco por el espacio academico para el
desarrollo de este trabajo.

% ----------------------------------------------------------
\section*{Referencias}
% ----------------------------------------------------------
\begingroup
\small
\begin{enumerate}
\setlength{\itemsep}{0pt}
\setlength{\parsep}{0pt}

\item F. Rosenblatt, ``The perceptron: A probabilistic model for
  information storage and organization in the brain,''
  \textit{Psychological Review}, vol. 65, no. 6, pp. 386--408, 1958.

\item M. Minsky and S. Papert, \textit{Perceptrons: An Introduction
  to Computational Geometry.} Cambridge, MA: MIT Press, 1969.

\item I. Goodfellow, Y. Bengio, and A. Courville,
  \textit{Deep Learning.} Cambridge, MA: MIT Press, 2016.

\item S. Haykin, \textit{Neural Networks and Learning Machines,}
  3rd ed. Upper Saddle River, NJ: Pearson, 2009.

\item C. M. Bishop, \textit{Pattern Recognition and Machine
  Learning.} New York: Springer, 2006.

\item V. Nair and G. E. Hinton, ``Rectified linear units improve
  restricted Boltzmann machines,'' in \textit{Proc. ICML},
  pp. 807--814, 2010.

\item X. Glorot, A. Bordes, and Y. Bengio, ``Deep sparse rectifier
  neural networks,'' in \textit{Proc. AISTATS},
  pp. 315--323, 2011.

\end{enumerate}
\endgroup

\end{document}
